Fix \(P\in\mathfrak P_{\mathcal D}^{\mathrm{obs}}\).  The family \(\mathcal D\) is finite and nonempty, and so is \(\mathcal D_{\mathrm{pred}}(P)\).  Positive costs, Assumption~\(\mathrm{ass{:}nondegenerate\text{-}components}\), and Theorem~\(\mathrm{thm{:}budget\text{-}profile}\) give, for every \(d\in\mathcal D\),
\[
 0<\mathcal V_d^*(P)
 =\left(\sqrt{c_{E,d}V_{E,d}}+
          \sqrt{c_{O,d}V_{O,d}}\right)^2<\infty.       \tag{1}
\]
Substitution of (1) into Definition~\(\mathrm{def{:}predictive\text{-}benchmark}\) proves the displayed identity for \(\rho_{\mathrm{pred}}(P)\).  Its numerator minimizes over a nonempty subset of the design set used by its denominator, so the ratio is at least one; (1) makes it finite.

For the best-tie diagnostic rule, a constant \(C\) is valid exactly when
\[
 \min_{d\in\mathcal D_{\mathrm{pred}}(P)}\mathcal V_d^*(P)
 \le C\min_{d'\in\mathcal D}\mathcal V_{d'}^*(P).
\]
The denominator is strictly positive by (1).  Dividing by it proves that the least valid constant is exactly \(\rho_{\mathrm{pred}}(P)\).

For the pairwise assertion, put
\[
 C_{\mathrm{pair}}(P)=
 \max_{\substack{d,d'\in\mathcal D:\,
 \Psi_{\mathrm{CLW},d}(P;B)\le
 \Psi_{\mathrm{CLW},d'}(P;B)}}
 \frac{\mathcal V_d^*(P)}{\mathcal V_{d'}^*(P)}.        \tag{2}
\]
The indexing set in (2) is nonempty because it contains every diagonal pair, and it is finite.  Every ratio in it is finite and positive by (1).  Hence \(C_{\mathrm{pair}}(P)<\infty\).  Any constant satisfying every diagnostic-order implication must dominate each ratio in (2), and hence their maximum.  Conversely, the maximum in (2) dominates every indexed ratio and therefore satisfies all the implications.  This proves exact minimality.

No causal restriction has been used in (1)--(2): outcome-bridge existence, selection-bridge existence, source-arm positivity, operator bijectivity, and nondegenerate score components are precisely the observable premises collected in \(\mathfrak P_{\mathcal D}^{\mathrm{obs}}\).  The remaining assertions concern the stronger binary causal subclass.  Theorem~\(\mathrm{thm{:}predictive\text{-}impossibility}\) proves for every fixed \(\chi>1\) that
\[
 \sup_{P\in\mathcal U_{\mathrm{bin}}}
 \rho_{\mathrm{pred}}(P)=\infty,
\]
while the uniquely diagnostic-selected center design has paired cost \(5\chi\) and the efficient center design has paired cost \(5\).  Lemma~\(\mathrm{lem{:}sharp\text{-}uniform\text{-}binary\text{-}elbow}\) proves that every sufficiently small-\(\epsilon\) minimizer over \(\chi>1\) satisfies
\[
 \chi_{\mathrm{opt}}(\epsilon)
 \sim\frac{\sqrt[3]{18}}{15}\epsilon^{-2/3},
\]
and proves the common arm-normalized and IKM-normalized expansion
\[
\begin{aligned}
 \epsilon^2\inf_{\chi>1}\rho_{\mathrm{pred}}(P_{\chi,\epsilon})
 &=\epsilon^2\inf_{\chi>1}
   \rho_{\mathrm{pred}}^{\mathrm{IKM}}(P_{\chi,\epsilon})\\
 &=\frac1{726}+\frac{5\sqrt[3]{12}}{968}\epsilon^{2/3}
   +o(\epsilon^{2/3}).
\end{aligned}
\]
Thus the finite-catalogue constant is exactly computable law by law on the observable bridge class, whereas the stated binary causal subclass prevents any finite primitive uniform guarantee for this particular diagnostic.  This proves the proposed rescoping and all remaining clauses.